\begin{tikzpicture}[scale=0.9]
\draw[rounded corners=5pt,fill=Mapcolor] (0,0)--++(5:16cm)--++(-95:3.5)--++(181:16.5cm)--cycle;
\draw[rounded corners=5pt,draw=white,line width=2pt,dashed,scale=0.9] (0.15,-0.15)--++(5:17.2cm)--++(-95:3.5)--++(181:17.7cm)--cycle;
\node[white,rotate=3] at (8,-0.6){\bf \Huge LỜI NÓI ĐẦU};
\end{tikzpicture}\par
Độc giả thân mến! 

Chúng tôi hiểu rằng việc ôn thi không chỉ đòi hỏi kiến thức sâu rộng mà còn yêu cầu phương pháp học tập và ôn luyện đúng đắn. Chính vì vậy, đây không chỉ là một quyển sách tham khảo mà còn là một người bạn dẫn dắt người học đi qua từng giai đoạn từ dễ đến khó để chinh phục được môn học tuy thú vị nhưng cũng rất khó khăn này.

Quyển sách này được biên soạn với mục tiêu giúp người học ôn luyện môn Vật Lí một cách hiệu quả và tự tin để chinh phục kì thi tốt nghiệp Trung học phổ thông.

Bộ sách được chia làm 4 phần: \textcolor{coolblack}{\textbf{NHIỆT - KHÍ - TỪ TRƯỜNG - HẠT NHÂN}}

Cấu trúc của quyển sách bao gồm: 
\begin{itemize}
	\item  \textit{Lí Thuyết Trọng Tâm}: được trình bày khoa học, kĩ lưỡng, đặc biệt là những hình ảnh minh họa phù hợp và đẹp mắt giúp người học nắm vững được nội dung chính của từng bài học.
	\item  \textit{Câu hỏi lí thuyết}: Giúp người học vận dụng ngay lập tức phần lí thuyết để trả lời các câu hỏi.
	\item \textit{Ví dụ minh họa}: Là những bài tập điển hình trong mỗi dạng toán đã được chia nhỏ. Giúp người học dễ dàng tiếp cận các dạng bài từ dễ đến khó.
	\item \textit{Bài tập về nhà}: Giúp người học tự rèn luyện thêm để làm chủ kiến thức.
	\item \textit{Bộ đề ôn luyện cuối chương}: Theo chuẩn cấu trúc đề thi của Bộ Giáo Dục, được sắp xếp với 5 đề đầu tiên cơ bản giúp các em hệ thống kiến thức cốt lõi và 10 đề nâng cao giúp các em chuẩn bị chinh phục điểm số $9^+$.
\end{itemize}

Sự nỗ lực của chúng tôi đã được thể hiện qua việc tìm kiếm, nghiên cứu và tổng hợp thông tin từ các nguồn đáng tin cậy, sưu tầm trên internet, chương trình giáo dục phổ thông 2018 đến các tài liệu chuẩn của Bộ Giáo dục và Đào tạo. 

Với sự cam kết và nỗ lực không ngừng, chúng tôi hy vọng rằng quyển sách này sẽ trở thành một nguồn tài liệu quan trọng, giúp người học ôn thi tốt nghiệp Trung học phổ thông môn Vật Lí một cách hiệu quả và tự tin. 

Trong quá trình biên soạn mặc dù đã cố gắng hết sức để đảm bảo sự chính xác và đầy đủ của mỗi phần nội dung. Tuy nhiên, chúng tôi hiểu rằng không có gì hoàn hảo và khó tránh khỏi những sai sót xuất hiện. Vì vậy rất mong nhận được sự thông cảm và phản hồi từ bạn đọc. 

Xin gửi lời cảm ơn sâu sắc đến tất cả những người đã đóng góp và hỗ trợ chúng tôi trong quá trình biên soạn sách này. 

Cuối cùng, xin chúc các em thành công và đạt được những kết quả tốt nhất trong hành trình học tập của mình!
\begin{flushright}
\parbox{0.8\textwidth}{
\begin{center}
Nhóm tác giả\\
\textcolor{coolblack}{\textbf{Vũ Ngọc Anh - TS. Trần Quốc Tuấn - Đinh Hoàng Tùng}}
\end{center}
}
\end{flushright}